% Заключение
\chapter*{ЗАКЛЮЧЕНИЕ}
\addcontentsline{toc}{chapter}{ЗАКЛЮЧЕНИЕ}

В ходе выполнения курсовой работы достигнута поставленная цель: спроектированы и реализованы база данных медицинской клиники и приложение доступа к ней, обеспечивающие корректную запись на прием, проведение приема, формирование счетов, учет оплат и кэширование часто используемых данных.

При выполнении работы получены следующие основные результаты:
\begin{enumerate}
    \item проведен анализ предметной области медицинской клиники и формализована задача разработки базы данных и приложения;
    \item спроектированы сущности базы данных, связи между ними, ограничения целостности, ролевая модель, процедуры и триггеры;
    \item выбраны средства реализации, соответствующие требованиям задания: PostgreSQL, Redis, Go, Gin, React и TypeScript;
    \item реализованы таблицы, индексы, домены, перечислимые типы, хранимые процедуры, триггеры и прикладной REST API;
    \item разработаны тестовые сценарии для проверки записи, отмены, проведения приема, формирования счетов и регистрации оплаты;
    \item исследованы конкурентный доступ и влияние индексов на производительность ключевых операций.
\end{enumerate}

Поставленные задачи решены в полном объеме: разработаны структура базы данных, ролевая модель, функции, триггеры, REST API, frontend-интерфейс и внешнее кэширование с использованием Redis.

Практическая проверка показала, что система запрещает конфликтующие записи по времени, корректно выполняет основные бизнес-операции и ускоряет чтение расписания и прайс-листа за счет кэширования. Результаты исследования подтвердили приемлемое время выполнения ключевых операций на выбранных объемах данных.

Разработанная система может использоваться как основа информационной системы небольшого медицинского центра. Для дальнейшего развития рекомендуется добавить журнал аудита действий пользователей, расширить управление расписанием нескольких врачей и реализовать резервное копирование базы данных.

\clearpage
